% Práctica 05 | Premios y canjes.
\section{Dominios y restricciones de premios y canjes}
Se definen las tablas \texttt{Premio}, \texttt{Canje} y \texttt{Disponer}
a partir del modelo relacional. Los identificadores conservan los nombres
y acentos mediante comillas dobles. La creación de tablas establece las
columnas y sus tipos; las restricciones se agregan con \texttt{ALTER TABLE}.

Todas las columnas son obligatorias (\texttt{NOT NULL}) y carecen de valor
predeterminado. Los enteros admiten un máximo de 2\,147\,483\,647; las
restricciones siguientes precisan su límite inferior. Las claves primarias
identifican las filas y las claves foráneas mantienen la integridad
referencial \citep{pg-constraints}.

\subsection{Premio}
El catálogo almacena el valor estimado, el costo en puntos y las
características de cada premio. La clave primaria es \texttt{idPremio}.

\begin{tabladominios}
\texttt{idPremio} & \texttt{integer} & Entero positivo; clave primaria. \\
\texttt{valorAproximado} & \texttt{integer} & Entero no negativo; importe estimado en pesos. Se admite cero. \\
\texttt{puntosNecesarios} & \texttt{integer} & Entero positivo; puntos requeridos por unidad. \\
\texttt{nombre} & \texttt{varchar(50)} & Hasta 50 caracteres; debe contener al menos un carácter que no sea blanco. \\
\texttt{categoría} & \texttt{varchar(50)} & Hasta 50 caracteres; debe contener al menos un carácter que no sea blanco. No se impone un catálogo cerrado. \\
\texttt{rangoEdad} & \texttt{varchar(50)} & Descriptor de hasta 50 caracteres; debe contener al menos un carácter que no sea blanco. \\
\end{tabladominios}

No se exige unicidad en \texttt{nombre}, pues distintos premios pueden
compartir una denominación. El inventario se registra por sucursal en
\texttt{Disponer}, de forma independiente del catálogo.

\clearpage
\subsection{Canje}
Cada registro identifica una operación de entrega de premios. La clave
primaria es \texttt{idCanje}; las referencias a cliente, premio, sucursal
y encargado pueden repetirse entre operaciones distintas.

\begin{tabladominios}
\texttt{idCanje} & \texttt{integer} & Entero positivo; clave primaria global del canje. \\
\texttt{idSucursal} & \texttt{integer} & Entero positivo; clave foránea.\newline\texttt{SUCURSAL(idSucursal)}. \\
\texttt{idCliente} & \texttt{integer} & Entero positivo; clave foránea.\newline\texttt{Cliente(idCliente)}. \\
\texttt{idPremio} & \texttt{integer} & Entero positivo; clave foránea.\newline\texttt{Premio(idPremio)}. \\
\texttt{CURP} & \texttt{char(18)} & Cadena de 18 caracteres con formato de CURP; referencia a la clave del encargado de premios. \\
\texttt{hora} & \texttt{time} & Hora válida sin zona horaria. \\
\texttt{puntosUtilizados} & \texttt{integer} & Entero positivo; puntos consumidos en la operación. \\
\texttt{cantidadPremios} & \texttt{integer} & Entero positivo; número de unidades entregadas. \\
\texttt{fecha} & \texttt{date} & Fecha válida y finita; se rechazan los valores infinito positivo e infinito negativo. \\
\end{tabladominios}

La clave foránea \texttt{CURP} referencia a
\texttt{Encargado de premios(CURP)}. La comprobación de formato exige
cuatro letras mayúsculas, seis dígitos, H o M, cinco letras mayúsculas
y dos caracteres alfanuméricos finales.

Esta comprobación valida la estructura de la cadena; la expedición oficial
de la CURP y la validez de su fecha interna quedan fuera de su alcance.
La fecha y la hora de la operación se almacenan en columnas independientes.

\clearpage
\subsection{Disponer}
Esta relación representa las existencias de cada premio en una sucursal.
La clave primaria compuesta \texttt{(idSucursal, idPremio)} identifica una
única fila de inventario por combinación de sucursal y premio.

\begin{tabladominios}
\texttt{idSucursal} & \texttt{integer} & Entero positivo; clave foránea.\newline\texttt{SUCURSAL(idSucursal)}. \\
\texttt{idPremio} & \texttt{integer} & Entero positivo; clave foránea.\newline\texttt{Premio(idPremio)}. \\
\texttt{cantidadDisponible} & \texttt{integer} & Entero no negativo; cero representa inventario agotado. \\
\end{tabladominios}

Un premio puede estar disponible en varias sucursales y cada sucursal puede
ofrecer varios premios. La clave compuesta evita duplicar una misma
combinación y el dominio impide almacenar existencias negativas.

\subsection{Integridad referencial}
Las seis claves foráneas utilizan \texttt{ON UPDATE NO ACTION} y
\texttt{ON DELETE NO ACTION}. Una clave referenciada no puede modificarse
o eliminarse si permanecen filas que dependen de ella. El historial de
canjes se conserva sin eliminaciones en cascada.

La suficiencia de puntos, la disponibilidad al realizar un canje, el
descuento de existencias y la relación entre el costo vigente y los puntos
consumidos requieren operaciones transaccionales entre tablas. Lo mismo
ocurre con la asignación del encargado a la sucursal. Estas reglas no se
atribuyen a las restricciones \texttt{CHECK} de cada fila.

\subsection{Validación de las restricciones}
Se ejecutaron 58 casos sobre un esquema de prueba aislado en PostgreSQL 18
y se repitieron satisfactoriamente en PostgreSQL 16.15. Se comprobaron datos
válidos, límites con cero, valores nulos en cada columna, dominios
inválidos, longitudes excesivas, fechas infinitas, claves primarias
duplicadas, las seis claves foráneas y la protección de las referencias.

Los rechazos esperados se comprobaron mediante SQLSTATE y, cuando
correspondía, por el nombre de la restricción. Los 58 casos fueron correctos.
La ejecución terminó con \texttt{ROLLBACK}, por lo que no dejó datos de
prueba. Las tablas externas se representaron mediante dependencias mínimas
para evaluar exclusivamente estas tres relaciones.
